\section{Introduction}
\label{sec:introduction}

Indonesian, the official language of the world's fourth-most populous country, has roughly 270 million speakers but a comparatively thin natural-language-processing infrastructure. The clitic \id{-nya} appears in approximately 19.5\% of suffix occurrences in Indonesian corpus data \citep{denistia2021morphology}, making it one of the highest-frequency morphological items in the language. It is also one of the most polysemous: a single surface form encodes at least four distinct grammatical functions, including third-person singular possessive (\id{rumahnya} `her house'), definite marker (\id{rumahnya} `the house', given prior reference), anaphoric pronoun referring to a previously mentioned entity (\id{pidatonya} `his/her speech'), and nominaliser (\id{datangnya} `the arrival').

Standard Indonesian information-retrieval pipelines treat \id{-nya} as a generic suffix to be stripped by a stemmer such as the Nazief--Adriani algorithm \citep{nazief1996indonesian} or its modern descendant Sastrawi \citep{sastrawi}. The most authoritative empirical study on Indonesian stemming and retrieval effectiveness \citep{tala2003stemming} predates BERT and dense retrieval entirely. The preprocessing defaults inherited from that pre-neural era propagate into modern retrieval systems built on multilingual dense encoders such as BGE-m3 \citep{chen2024bgem3} and multilingual E5 \citep{wang2024multilingual_e5}, where they have not been controlled-tested. Recent surveys of preprocessing on Transformer-based models report that classical preprocessing often does not help, or actively hurts, modern neural systems on English-centric benchmarks \citep{camacho2023preprocessing}. Whether that finding generalises to a morphologically distinctive low-resource language with a high-frequency polysemous clitic is an open empirical question.

A separate and equally untested concern is the \emph{quality} of the \id{-nya} stripping itself. A naive regex of the form \texttt{(\textbackslash w+?)nya\textbackslash b} --- which appears in widely-distributed Indonesian NLP tutorials --- over-strips a substantial set of high-frequency Indonesian non-clitic words such as \id{punya} `to have', \id{tanya} `to ask', \id{hanya} `only', and \id{biasanya} `usually', as well as loanwords and proper nouns including \id{Kenya}, \id{Sonya}, and the personal name \id{Tanya}. Quantifying the harm caused by this naive default is a public-good contribution distinct from the more sophisticated question of whether \emph{any} preprocessing helps modern retrievers.

This paper isolates Indonesian \id{-nya} preprocessing as an experimental variable in modern retrieval. We address three primary questions:

\begin{enumerate}
    \item \textbf{What works best?} Among five preprocessing strategies (keep, naive strip, dictionary-aware Sastrawi clitic strip, sentinel mask, oracle anaphora resolution), which yields the highest retrieval effectiveness on MIRACL-id?
    \item \textbf{Does the answer depend on retriever architecture?} Specifically, does the optimal preprocessing strategy differ between BM25 and BGE-m3?
    \item \textbf{What is the upper bound?} How large is the ceiling, established by oracle anaphora resolution, on the gain attainable from \id{-nya}-aware preprocessing?
\end{enumerate}

Our contributions are threefold:

\begin{itemize}
    \item The first controlled study isolating Indonesian \id{-nya} preprocessing as an experimental variable in a modern dense-retrieval setting.
    \item Empirical quantification of the harm caused by naive regex \id{-nya} stripping --- the harm-baseline that, to our knowledge, has not been measured despite the prevalence of the pattern in practitioner code.
    \item An oracle ceiling that bounds the value of future research into Indonesian anaphora resolution for retrieval applications.
\end{itemize}

The remainder of the paper is organised as follows. Section~\ref{sec:related} situates the work in three adjacent literatures: Indonesian morphology, Indonesian retrieval evaluation, and anaphora-aware retrieval. Section~\ref{sec:method} specifies the five preprocessing conditions and the experimental design. Section~\ref{sec:experiments} describes implementation details, and Section~\ref{sec:results} presents results. Section~\ref{sec:discussion} interprets the findings and discusses limitations. Section~\ref{sec:conclusion} concludes.
